\documentclass[reprint,aps,prb,superscriptaddress]{revtex4-2}

\usepackage{graphicx}
\usepackage{amsmath}
\usepackage{amssymb}
\usepackage{bm}
\usepackage{hyperref}
\usepackage[utf8]{inputenc}

\begin{document}

\title{GPU-Accelerated Tensor Network Dynamics of Radical Pair Spin Systems: \\ 
Breaking the $\chi=1500$ Barrier on Consumer Hardware}

\author{The Architect}
\affiliation{SUBSTRATE Research Lab}

\author{Antigravity AI}
\affiliation{Advanced Agentic Coding, Google DeepMind}

\date{\today}

\begin{abstract}
The simulation of radical pair spin dynamics in biological systems, such as the avian magnetic compass (Cryptochrome 4a), 
is limited by the exponential growth of the Hilbert space. Matrix Product State (MPS) methods have successfully 
extended these simulations, yet previous implementations were capped at bond dimensions of $\chi \le 1500$ due 
to CPU memory constraints. In this work, we present SUBSTRATE-GPU, an optimized solver utilizing NVIDIA Blackwell 
architectures and selective recomputation to achieve $\chi=2500$ for systems of up to 40 nuclear spins. 
Our benchmarks demonstrate a significant reduction in wall-time and VRAM management strategies that enable 
high-fidelity simulations previously restricted to supercomputing clusters.
\end{abstract}

\maketitle

\section{Introduction}
The radical pair mechanism (RPM) is the primary model for magnetoreception in migratory birds...

\section{Methods}
\subsection{Tensor Network Representation}
We represent the spin density matrix in Liouville space as a Matrix Product State (MPS)...

\subsection{GPU Optimization}
Our solver utilizes the \texttt{cuTensorNet} and \texttt{CuPy} libraries to offload tensor contractions 
to the GPU... We implement a selective recompute strategy for the environments to stay within 
the 16 GB VRAM limit of the RTX 5060 Ti at $\chi=2500$.

\section{Results}
\subsection{$\chi$-Convergence}
As shown in Fig. \ref{fig:scaling}, our implementation achieves convergence at bond dimensions...

\begin{figure}[t]
    \centering
    %\includegraphics[width=\linewidth]{chi_scaling.pdf}
    \caption{Scaling of wall-time per step and VRAM usage as a function of bond dimension $\chi$ for the 40-site ErCry4a system.}
    \label{fig:scaling}
\end{figure}

\section{Conclusion}
The ability to simulate $\chi=2500$ on consumer-grade hardware marks a milestone...

\begin{acknowledgments}
Supported by the SUBSTRATE Sovereign Infrastructure.
\end{acknowledgments}

\bibliography{refs}

\end{document}
